\documentclass[11pt]{article}
\usepackage[margin=1in]{geometry}
\usepackage{amsmath,amssymb,amsthm}
\usepackage[hidelinks]{hyperref}
\newtheorem{theorem}{Theorem}
\newtheorem{lemma}{Lemma}
\newcommand{\E}{\mathbb E}
\newcommand{\Pp}{\mathbb P}
\title{The Laplace-transform threshold for \texorpdfstring{$(1+mse^s)^{-1}$}{(1+m s exp(s)) inverse}}
\author{MathDB problem 350435}
\date{12 September 2026}
\begin{document}
\maketitle

\section*{Task 1. Solve the Laplace-transform criterion}

Kesidis, Konstantopoulos and Zazanis conjectured the following classification
in Conjecture 1 of \cite[p.~18]{KKZ}. Their probabilistic construction already
established the sufficient range. We prove the missing necessity and give
a direct construction for sufficiency.

\begin{theorem}
For a real parameter $m\geq0$, the function
\[
 L_m(s)=\frac{1}{1+mse^s},\qquad s\geq0,
\]
is the Laplace transform of a probability distribution if and only if
$m=0$ or $m\geq e$. The conclusion holds even if the putative random
variable is initially allowed to take negative values.
\end{theorem}

We use the following elementary form of the nonnegative-coefficient
singularity principle, often called Pringsheim's theorem. Its proof
is included to make the argument self-contained.

\begin{lemma}\label{lem:pringsheim}
Let $f(z)=\sum_{n\geq0}a_nz^n$, where $a_n\geq0$, have radius of
convergence $R$ with $0<R<\infty$. Then $f$ has no holomorphic
continuation to a neighborhood of $R$ agreeing with the series just
to the left of $R$.
\end{lemma}
\begin{proof}
Suppose such a continuation exists on $|z-R|<\delta$, where
$0<\delta<R$. Set $b=R-\delta/4$ and $h=\delta/3$.
The disc $|z-b|<\delta/2$ is contained in $|z-R|<\delta$.
Consequently the Taylor series at $b$ converges at $b+h$.
Since $0<b<R$, termwise differentiation of the original series gives
\[
 \frac{f^{(k)}(b)}{k!}
   =\sum_{n\geq k}a_n\binom{n}{k}b^{n-k}\geq0.
\]
All terms in the next double series are nonnegative, so summation
can be interchanged:
\[
 \infty>
 \sum_{k\geq0}\frac{f^{(k)}(b)}{k!}h^k
 =\sum_{n\geq0}a_n\sum_{k=0}^n\binom{n}{k}b^{n-k}h^k
 =\sum_{n\geq0}a_n(b+h)^n .
\]
But $b+h=R+\delta/12>R$, contradicting the radius of convergence.
\end{proof}

\begin{proof}[Proof of the theorem: necessity]
Suppose $0<m<e$ and $L_m(s)=\E e^{-sX}$ for every $s\geq0$.
As $s\to\infty$, $L_m(s)\to0$. If $\Pp(X<0)>0$, then
$\Pp(X\leq-\varepsilon)>0$ for some $\varepsilon>0$, and hence
$\E e^{-sX}\geq e^{s\varepsilon}\Pp(X\leq-\varepsilon)\to\infty$.
Thus $X\geq0$ almost surely.

Differentiating its Laplace transform at $s>0$ gives
\[
 (-1)^nL_m^{(n)}(s)=\E[X^ne^{-sX}]\geq0.
\]
Differentiation is justified locally by boundedness of
$x^ne^{-cx}$ on $x\geq0$ for each $c>0$.
The explicit expression for $L_m$ is analytic at zero, so its
derivatives extend continuously there. It follows that
\[
 F(z):=L_m(-z)=\frac{1}{1-mze^{-z}}
       =\sum_{n\geq0}a_nz^n
 \quad\hbox{near }0,\qquad
 a_n=\frac{(-1)^nL_m^{(n)}(0)}{n!}\geq0.
 \tag{1}
\]
In particular, $a_0=1$ and $a_1=m>0$.

For every real $t\geq0$, the elementary inequality $te^{-t}\leq1/e$
shows that
\[
 1-mte^{-t}\geq1-m/e>0,\qquad
 0<F(t)\leq\frac{1}{1-m/e}.                 \tag{2}
\]
The displayed quotient therefore provides a holomorphic continuation
of the germ in (1) to a neighborhood of every point of the positive
real axis. If the Taylor radius in (1) were finite, its positive
endpoint would contradict Lemma~\ref{lem:pringsheim}. Thus the radius
is infinite. The series then equals the quotient on the positive
axis by analytic continuation. Nonnegativity of its coefficients
implies
\[
 F(t)=\sum_{n\geq0}a_nt^n\geq1+mt\qquad(t\geq0),
\]
which contradicts the uniform bound in (2). This excludes the whole
interval $0<m<e$.
\end{proof}

\begin{proof}[Proof of the theorem: sufficiency]
For $m=0$, take $X=0$ almost surely. If $m\geq e$, choose
$\lambda\in(0,1]$ so that
\[
 m=\frac{e^\lambda}{\lambda}.
\]
Such a choice exists because $e^\lambda/\lambda$ is continuous,
decreases from infinity to $e$ on $(0,1]$, and equals $e$ at $1$.
Put $p=e^{-\lambda}$. Let $Z$ be exponential with rate $\lambda$,
let $J$ have distribution
\[
 \Pp(J=j)=p(1-p)^j,\qquad j=0,1,\ldots,
\]
and let $Y_1,Y_2,\ldots$ have common density
\[
 q(y)=\frac{\lambda e^{-\lambda y}}{1-p}
             \boldsymbol{1}_{(0,1)}(y).
\]
Take all these variables independent and define
$T=1+Z+\sum_{i=1}^JY_i$. This is a finite nonnegative random
variable almost surely. The common transform of the $Y_i$ is
\[
 \phi(s)=\frac{\lambda(1-e^{-(\lambda+s)})}
                    {(\lambda+s)(1-p)},\qquad s\geq0 .
\]
In particular, $(1-p)\phi(s)\leq1-p<1$, so the geometric series
for the transform of the random sum converges. Independence yields
\begin{align*}
 \E e^{-sT}
 &=e^{-s}\frac{\lambda}{\lambda+s}
           \frac{p}{1-(1-p)\phi(s)}\\
 &=\frac{\lambda e^{-(\lambda+s)}}
             {s+\lambda e^{-(\lambda+s)}}
  =\frac{1}{1+(e^\lambda/\lambda)se^s}
  =L_m(s).
\end{align*}
This includes the endpoint $m=e$, with $\lambda=1$.
\end{proof}

\paragraph{Scope.}
The argument establishes the complete parameter classification,
not just a finite collection of derivative-sign tests. In fact,
the necessity proof shows that for every $0<m<e$ at least one
Taylor coefficient of $L_m(-z)$ is negative. The sufficiency proof
constructs an absolutely continuous probability law supported in
$[1,\infty)$ for every $m\geq e$.

\section*{Task 2. Quantify the first failed derivative sign}
For $0<m<e$, determine or sharply bound
\[
 n_*(m)=\min\{n\geq0:[z^n](1-mze^{-z})^{-1}<0\}.
\]
In particular, obtain its asymptotic behavior as $m\uparrow e$.
The proof above guarantees finiteness but gives no effective bound.

\section*{Task 3. Locate the sign changes of the inverse transform}
For $0<m<e$, study the piecewise polynomial inverse Laplace transform
of $L_m$. Determine its first negative interval and obtain rigorous
large-time asymptotics, including the transition as $m\uparrow e$.

\section*{Task 4. Classify fractional-power variants}
For $m,\tau>0$ and $0<\alpha\leq1$, determine precisely when
\[
 (1+m s^\alpha e^{\tau s})^{-1}
\]
is a probability Laplace transform. Find stochastic constructions
throughout the admissible parameter region.

\section*{Task 5. Determine the shape of the admissible distributions}
For the probability laws $T_m$, $m\geq e$, determine whether their
densities are log-concave and classify the monotonicity of their hazard
rates. Obtain sharp bounds that remain valid at the critical value $m=e$.

\begin{thebibliography}{9}
\bibitem{KKZ}
G. Kesidis, T. Konstantopoulos and M. A. Zazanis,
\emph{Age of information without service preemption},
\href{https://arxiv.org/abs/2104.08050v3}{arXiv:2104.08050v3}
(2021), Conjecture 1, p.~18.
\end{thebibliography}
\end{document}
